\documentclass[10pt,twocolumn]{article}
\usepackage[T1]{fontenc}
\usepackage[utf8]{inputenc}
\usepackage{lmodern}
\usepackage[margin=1.8cm,top=2.4cm,bottom=2cm,columnsep=0.6cm]{geometry}
\usepackage{fancyhdr}
\usepackage{titlesec}
\usepackage{booktabs}
\usepackage{amsmath,amssymb,amsfonts}
\usepackage{microtype}
\usepackage{xcolor}
\usepackage{mdframed}
\usepackage{parskip}
\usepackage{enumitem}
\usepackage{hyperref}

\definecolor{rulered}{RGB}{180,30,30}
\definecolor{boxgray}{RGB}{245,245,245}
\definecolor{keywordgray}{RGB}{100,100,100}

\pagestyle{fancy}
\fancyhf{}
\fancyhead[L]{\textsc{\small Claude Mag}}
\fancyhead[C]{\small\textit{Machine Learning Research Digest}}
\fancyhead[R]{\small 2026-09-09}
\fancyfoot[C]{\small\thepage}
\renewcommand{\headrulewidth}{0.8pt}

\titleformat{\section}{\normalsize\bfseries\color{rulered}}{}{0pt}{}%
  [\vspace{-4pt}\color{rulered}\rule{\columnwidth}{0.4pt}]
\titlespacing{\section}{0pt}{8pt}{4pt}

\hypersetup{colorlinks=true,urlcolor=rulered,linkcolor=black}

\begin{document}
\setlength{\parindent}{0pt}
\setlength{\parskip}{4pt}

{\small\color{keywordgray}\textbf{Keywords:}
  diffusion transformer \textbullet{}
  image generation \textbullet{}
  open training recipe \textbullet{}
  Muon optimizer}\par\vspace{4pt}

{\LARGE\bfseries\raggedright
  LLaDA-Image: Building Strong Image Generators
  with Fully Open Training Recipes}\par\vspace{4pt}

{\large\itshape\raggedright
  A 6B Diffusion Transformer with Staged Training,
  Parameter-Free RMSNorm, and the Muon Optimizer
  Achieves State-of-the-Art on Qwen-Image-Bench
  with Full Open Release}\par\vspace{6pt}

{\small\textbf{By} Chuyan Chen, Haoxing Chen, et al.\
  (inclusionAI, Ant Group)}\par\vspace{2pt}

{\small\color{keywordgray}arXiv:2609.03796 \textbullet{} Preprint, September 2026}
\par\vspace{2pt}\noindent{\color{rulered}\rule{\columnwidth}{0.6pt}}\vspace{6pt}

High-quality text-to-image generation has long been dominated by
closed proprietary systems. LLaDA-Image breaks this pattern: a 6B
Diffusion Transformer (DiT) trained from scratch within a unified
framework alongside a frozen vision-language understanding module,
achieving state-of-the-art scores of 53.53 (English) and 53.38
(Chinese) on Qwen-Image-Bench, with all model weights, training code,
and stage-by-stage recipes publicly released. The key engineering
decisions --- staged training (image-only first, then paired text-image),
parameter-free RMSNorm, and the Muon optimizer --- are each responsible
for significant quality and stability gains and are now fully reproducible
by the community.

\begin{mdframed}[backgroundcolor=boxgray,linecolor=rulered,linewidth=1pt,
    innerleftmargin=6pt,innerrightmargin=6pt,
    innertopmargin=6pt,innerbottommargin=6pt]
\textbf{\small AT A GLANCE}\par\vspace{4pt}
\begin{description}[leftmargin=0pt,labelsep=4pt,itemsep=2pt,parsep=0pt,
    font=\small\bfseries]
  \item[Problem:] Competitive text-to-image generation requires
    proprietary recipes, closed weights, and unreported training
    decisions.
  \item[Why it matters:] Without reproducible open baselines, the
    community cannot build on, audit, or improve SOTA image generators.
  \item[Method:] 6B DiT + frozen LLaDA2.0-Mini backbone; image-only
    pre-training then paired text-image training; parameter-free RMSNorm;
    Muon optimizer.
  \item[Result:] SOTA Qwen-Image-Bench 53.53 EN / 53.38 ZH; Muon
    reduces training compute by 20\%; Turbo variant in 2--4 steps.
  \item[Evidence:] Qwen-Image-Bench, GenEval, T2I-CompBench; ablations
    on staged training, RMSNorm, and optimizer.
\end{description}
\end{mdframed}

\section{Architecture: Unified DiT Plus Frozen LLM}

LLaDA-Image pairs two modules:

\textbf{Image Generation Module (DiT).} A 6-billion-parameter Diffusion
Transformer trained from scratch to perform masked discrete diffusion
over image tokens. The DiT models the conditional distribution
$p(x_{\text{image}} | c_{\text{text}})$ where $c_{\text{text}}$ is the
conditioning signal from the understanding module.

\textbf{Vision-Language Understanding Module.} A frozen LLaDA2.0-Mini
diffusion language model, which processes text prompts and provides
conditioning for the DiT. Because LLaDA2.0-Mini was trained on both
English and Chinese text, the system natively understands Chinese
generation instructions --- a key advantage over English-dominant
baselines.

The two modules share an image tokeniser (VQ-VAE) that discretises
images into sequences of tokens, enabling a unified token vocabulary.

\section{Key Engineering Decision 1: Staged Training}

Training starts with image-only pre-training before any text-image
pair is introduced.

\textbf{Stage 1 --- Image-only pre-training (98M samples):}
The DiT is trained on unlabeled images using the masked diffusion
objective with no conditioning signal. This establishes a strong
visual generative prior: the model learns the manifold of natural
images before it is asked to condition on language.

\textbf{Stage 2 --- Paired text-image mid-training (122M samples):}
The DiT is fine-tuned on (text, image) pairs with conditioning from
the understanding module. Starting from a strong visual prior produces
significantly better text-following than training on pairs from scratch.

\textbf{Stage 3 --- Instruction fine-tuning:}
A smaller high-quality stage sharpens precision on detailed captions
and editing instructions.

The ablation is decisive: removing Stage 1 and starting with paired
training reduces Qwen-Image-Bench score by 3.8 points.

\section{Key Engineering Decision 2: Parameter-Free RMSNorm}

Standard affine RMSNorm applies learned scale $\gamma$ and bias
$\beta$:
\[
  \text{RMSNorm}_\text{affine}(x) = \gamma \cdot \frac{x}{\|x\|_2} + \beta
\]
At 6B parameters, the learned scale can drift during training, causing
loss spikes that require manual intervention.

Parameter-free RMSNorm simply normalises to unit norm:
\[
  \text{RMSNorm}_\text{free}(x) = \frac{x}{\|x\|_2}
\]
Fixing the output norm eliminates the scale drift channel, enabling
stable training throughout. Applied uniformly throughout the DiT,
this removes approximately 12M parameters while improving training
stability.

\section{Key Engineering Decision 3: Muon Optimizer}

Standard Adam-family optimisers update parameters along the raw
gradient direction, which can inflate or collapse parameter norms.
The Muon optimizer applies orthogonalised gradient updates:
\[
  \mathbf{G}_t^\perp = \operatorname{NewtonSchulz5}(\mathbf{M}_t)
\]
\[
  \theta_{t+1} = \theta_t - \eta \mathbf{G}_t^\perp
\]
where $\mathbf{M}_t$ is the momentum-corrected gradient and
$\operatorname{NewtonSchulz5}$ applies five iterations of the
Newton-Schulz iteration to orthogonalise $\mathbf{M}_t$. Orthogonal
updates preserve parameter norms, improving both convergence speed
and final quality.

\section{Masked Diffusion Objective}

For an image token sequence $x = (x_1,\ldots,x_T)$, the DiT is
trained to predict masked tokens:
\[
  \mathcal{L}_\text{DiT} =
  -\mathbb{E}_{t,\mathbf{m}}\!
  \left[ \sum_{i: m_i=1} \log p_\theta(x_i | x_{\setminus m}, c, t) \right]
\]
where $\mathbf{m}$ is a random mask, $x_{\setminus m}$ are unmasked
tokens, $c$ is text conditioning, and $t$ is the diffusion timestep.

\section{LLaDA-Image-Turbo: Distillation to 2--4 Steps}

LLaDA-Image-Turbo is distilled from the full model using consistency
distillation:
\[
  \mathcal{L}_\text{turbo} =
  \mathbb{E}_{x, t, s}
  \bigl\| f_\theta(x_t, t) - f_{\theta^-}(x_s, s) \bigr\|^2
\]
where $t > s$ are timesteps on the same trajectory, $f_\theta$ is the
student, and $f_{\theta^-}$ is an EMA teacher. Enforcing consistency
along generation trajectories enables convergence in 2--4 steps at
94--97\% of full model quality.

\section{Experimental Results}

\begin{center}
\begin{tabular}{lccc}
\toprule
Model & QB-EN & QB-ZH & GenEval \\
\midrule
DALL-E 3         & 50.1 & 48.7 & 0.61 \\
SD 3.5-L         & 49.6 & ---  & 0.71 \\
Flux-1.1-pro     & 51.4 & ---  & 0.72 \\
\textbf{LLaDA-Image} & \textbf{53.53} & \textbf{53.38} & \textbf{0.74} \\
LLaDA-Turbo (4-step) & 51.8 & 51.4 & 0.71 \\
\bottomrule
\end{tabular}
\end{center}

\textbf{Training compute:} Muon reduces GPU-hours to convergence by
20\% (75,500 vs.\ 94,400 with AdamW) and reaches the same validation
loss 25\% fewer steps.

\section{Ablations}

\textbf{Stage 1 removal:} $-3.8$ Qwen-Image-Bench points.

\textbf{Affine vs.\ free RMSNorm:} Affine causes training instability
at 4B$+$ scale; free RMSNorm eliminates all spikes.

\textbf{Muon vs.\ AdamW:} Muon is $+1.1$ Qwen-Image-Bench points
at the same compute budget.

\textbf{Turbo step count:} 4 steps retains 97\% of quality; 2 steps
94\%, at 8$\times$ speedup.

\textbf{Chinese prompt following:} The bilingual LLaDA2.0-Mini backbone
provides SOTA Chinese scores where English-dominant baselines degrade
significantly.

\section*{Reference}

Chuyan Chen et al.\ (inclusionAI).
\textbf{LLaDA-Image: Building Strong Image Generators with Fully Open
Training Recipes.}
arXiv:2609.03796, September 2026.
\url{https://arxiv.org/abs/2609.03796}

\end{document}
